\section{Completabilidad antes de tomar el logaritmo}
Sea $A$ la máscara de casillas legales contra las piezas anteriores y $j$ su tamaño.
Un punto de torres bistocástico soportado en $A$ garantiza $\operatorname{per}A>0$:
para cualquier conjunto de filas, su masa exige al menos tantas columnas vecinas.
Hall da el emparejamiento. Denotemos por $\sigma$ un emparejamiento uniforme en $A$.
Para $U$ un conjunto de sitios se usa $p_U=\Pr_{\rm torres}(U\subseteq\sigma)$,
incluido el valor cero si es incompatible. Esta ley sólo requiere el permanente positivo.

\paragraph{Entrada exacta y su dominio.}
En el grafo de ataques diagonales sea $I(T)$ el indicador de que $T$ no contiene ataques
y $a_U=\sum_{T\subseteq U}(-1)^{|U|-|T|}I(T)$.
En el álgebra booleana de activaciones, con $\rho\ge2$, pongamos
\[
 Z(z)=\sum_U a_U p_U z^U,\qquad
 \|f\|_\rho=\sum_U\rho^{|U|}|[U]f|.
\]
Por inversión de Möbius, para cada activación real $Y$,
$Z(\mathbf1_Y)=\mathbb E_{\rm torres} I(Y\cap\sigma)$. En particular
$Z(\mathbf1)=Q(B)/\operatorname{per}A$ incluso si $Q(B)=0$.

Fijemos los parámetros uniformes de difusión, brecha y aproximación de menores
de \emph{Media franjas}, partes 01--06. Su prueba de la parte 06, antes del logaritmo,
construye una unidad $D_h$ positiva en toda evaluación real. Para $j$ grande y
$\rho^2(\epsilon_A+1/j)\le\eta$ se obtiene $\|D_h^{-1}Z-1\|_\rho<1/2$.
En esta fuente los datos del gas y los pesos de supresión se forman con $p_U$,
el escalado de torres y sus menores. El logaritmo auxiliar $G(h)$ es del gas basal
y $D_h= \exp(G(h)+L(m)-h[m])$ se evalúa en campos reales.
La recuperación se hace para $Z$ definido arriba, antes de introducir $\log Z$.
El suceso histórico de éxito usado para medias posteriores queda fuera de este paso.

\paragraph{Implicación revisada.}
Evaluar en $\mathbf1_Y$ es un homomorfismo de norma a lo sumo uno. Luego
\[
 \frac{Z(\mathbf1_Y)}{D_h(\mathbf1_Y)}>\frac12,\qquad
                       Q(B)>0\quad(Y\text{ total}).
\]
$D_h$ puede ser pequeño al crecer $j$. Para esta deducción sólo importa su signo.
Sea $L_{\max}$ la mayor longitud disponible de una diagonal y
$j\max P_e\le M_0$ para el punto de torres $P$ de la máscara.
Entonces $\epsilon_A\le M_0L_{\max}/j$. Fijemos $\alpha>0$ con
$M_0\alpha<\min(\epsilon_{\mathcal D},\eta/(2\rho^2))$.
Por integralidad,
$L_{\max}<\lceil\alpha j\rceil$ implica $L_{\max}<\alpha j$.
En toda llegada histórica regular de estas características se cumple la hipótesis
de carga para $j$ grande y existe una completación. Es una afirmación determinista
condicionada a las estimaciones de torres citadas, sobre tableros con esa carga.
